

\section{L'impasse de l'enseignement explicite : Saturation déclarative, négligence procédurale}
space{0.3cm}
oindent

L'histoire de l'enseignement des langues secondes (L2) est marquée par une tension persistante entre
la théorie intuitive de l'apprentissage scolaire et la réalité biologique de l'acquisition du
langage. Pendant des siècles, et de manière prédominante dans les systèmes éducatifs contemporains,
l'approche pédagogique standard a reposé sur un postulat rationaliste : la langue est un système de
règles logiques qui, une fois comprises intellectuellement, mémorisées consciemment et pratiquées
délibérément, se transforment en une compétence fluide. Cette approche, souvent incarnée par la
méthode \og Grammaire-Traduction\fg{} (GTM) ou, plus récemment, par des formes rigides
d'enseignement explicite (Focus on Forms), place l'analyse métalinguistique au cœur du processus
d'apprentissage.

Cependant, un fossé grandissant sépare les résultats de cette méthode des objectifs de communication
réelle. Les enseignants et les apprenants font face à un phénomène récurrent et frustrant : l'élève
capable de réciter des règles complexes de conjugaison et de réussir des exercices à trous avec une
précision chirurgicale se trouve souvent incapable de formuler une phrase cohérente en temps réel
lors d'une conversation spontanée. Ce décalage, loin d'être un échec individuel ou un manque de \og
talent\fg{}, est symptomatique d'une erreur catégorielle fondamentale dans la conception même de
l'apprentissage linguistique.

La recherche en neurosciences cognitives, et plus particulièrement le Modèle Déclaratif/Procédural
(DP) théorisé par Michael Ullman, suggère que l'enseignement explicite traditionnel conduit les
apprenants dans une véritable impasse neurobiologique. En focalisant l'effort cognitif sur la
mémorisation de règles et le traitement conscient (déclaratif), la pédagogie classique sature les
ressources attentionnelles limitées de l'apprenant tout en échouant à stimuler les circuits
neuronaux responsables de l'automatisation (procédurale). Pire encore, des preuves émergentes
indiquent que l'hyper-sollicitation du système déclaratif pourrait activement inhiber le
développement de la compétence procédurale, un phénomène connu sous le nom d'effet de bascule (\og
see-saw effect\fg{}).

Il ne s'agit pas simplement de comparer
des méthodes pédagogiques, mais d'explorer l'architecture de la mémoire humaine pour comprendre
pourquoi la voie de l'enseignement explicite, bien que séduisante par sa clarté apparente, se heurte
aux limites physiologiques du cerveau.

Nous articulerons notre analyse autour de quatre axes majeurs. Premièrement, nous détaillerons
l'anatomie fonctionnelle du modèle Déclaratif/Procédural pour établir la distinction cruciale entre
le \og savoir que\fg{} (lexique) et le \og savoir comment\fg{} (grammaire). Deuxièmement, nous
examinerons la dynamique de compétition neuronale entre l'hippocampe et les ganglions de la base,
démontrant comment l'enseignement explicite peut agir comme un frein neurologique à la fluidité.
Troisièmement, nous intégrerons la Théorie de la Charge Cognitive de John Sweller pour illustrer
comment l'analyse grammaticale et la traduction mentale saturent la mémoire de travail, créant un
goulot d'étranglement infranchissable. Enfin, nous déconstruirons le mythe de la \og conversion\fg{}
des connaissances (Skill Acquisition Theory), en mettant en lumière l'impossibilité mathématique et
temporelle de procéduraliser des règles explicites dans le cadre restreint de la salle de classe
scolaire.

Pour saisir la nature de l'impasse pédagogique actuelle, il est impératif de comprendre que le
langage, du point de vue cérébral, n'est pas une entité monolithique. Contrairement aux théories
linguistiques modulaires qui postulaient l'existence d'un \og organe du langage\fg{} dédié, les
neurosciences modernes, appuyées par le modèle DP, révèlent que le langage \og coopte\fg{} deux
systèmes de mémoire à long terme distincts, ayant évolué bien avant l'apparition de la parole : la
mémoire déclarative et la mémoire procédurale.



\subsection{La mémoire déclarative : le siège du lexique et des faits}

La mémoire déclarative est le système neurocognitif dédié à l'apprentissage, au stockage et à la
récupération des connaissances explicites sur le monde (mémoire sémantique) et des expériences
personnelles (mémoire épisodique). Dans le contexte linguistique, ce système est le substrat du
\textbf{lexique mental}.



\subsection{Caractéristiques fonctionnelles et cognitives}

Le système déclaratif possède des propriétés uniques qui le rendent particulièrement adapté à
l'apprentissage du vocabulaire, mais inadapté à la gestion de la grammaire en temps réel.

\begin{itemize}\item \textbf{Apprentissage Rapide (\og One-Shot Learning\fg{}) :} La mémoire déclarative est capable d'encoder de nouvelles associations arbitraires après une seule exposition. C'est ce mécanisme qui permet d'apprendre que le son \og chat\fg{} correspond à l'animal félin, une association qui ne suit aucune règle logique déductible.\cite{II-2-2-ref1}\item \textbf{Accès Conscient et Explicite :} Les connaissances stockées dans ce système sont accessibles à la conscience. L'apprenant \og sait qu'il sait\fg{} et peut verbaliser cette connaissance (ex: \og Le mot pour 'livre' en anglais est 'book'\fg{}).\item \textbf{Flexibilité et Associativité :} Ce système est relationnel ; il permet de lier des informations disparates entre elles (contexte, sens, son), mais il n'est pas spécialisé dans le traitement séquentiel rapide.\cite{II-2-2-ref2}\item \textbf{Dépendance Attentionnelle :} L'encodage et la récupération en mémoire déclarative exigent des ressources attentionnelles importantes. C'est un système coûteux en énergie cognitive.\end{itemize}

\subsection{Substrats neuroanatomiques}

Le siège anatomique de la mémoire déclarative se situe dans le lobe temporal médial (MTL), avec une
importance critique de l'hippocampe pour l'acquisition initiale et la consolidation des souvenirs.
Les régions néocorticales temporales (comme le gyrus temporal supérieur) et pariétales (notamment le
gyrus supramarginal pour le stockage phonologique) jouent également un rôle clé dans le stockage à
long terme des connaissances lexicales.\cite{II-2-2-ref2}



\subsection{Le rôle dans l'apprentissage l2}

Chez l'adulte apprenant une langue seconde, la mémoire déclarative tend à prendre en charge des
fonctions qui, chez le locuteur natif, sont gérées par le système procédural. En raison de sa
plasticité qui persiste tout au long de la vie (contrairement à la mémoire procédurale qui pourrait
décliner ou devenir moins accessible après la puberté), la mémoire déclarative devient le système
\og par défaut\fg{} pour l'apprenant tardif. L'apprenant stocke non seulement les mots, mais aussi
des \og chunks\fg{} de phrases et des règles grammaticales explicites sous forme de faits
déclaratifs (ex: \og Il faut ajouter \-ed pour le passé\fg{}).\cite{II-2-2-ref1}



\subsection{La mémoire procédurale : le moteur de la grammaire et des automatismes}

À l'opposé du spectre cognitif se trouve la mémoire procédurale, un système implicite qui sous-tend
l'apprentissage et le contrôle des habiletés motrices et cognitives. Dans le cadre du modèle DP, ce
système est le substrat de la \textbf{grammaire mentale} (syntaxe, morphologie régulière,
phonologie).



\subsection{Caractéristiques fonctionnelles et cognitives}

La mémoire procédurale fonctionne selon des principes radicalement différents de ceux de la mémoire
déclarative, ce qui explique pourquoi l'enseignement explicite (basé sur le déclaratif) échoue
souvent à l'activer.

\begin{itemize}\item \textbf{Apprentissage Graduel et Statistique :} Contrairement à l'apprentissage \og one-shot\fg{} du déclaratif, le système procédural apprend lentement. Il nécessite une exposition répétée et massive à des stimuli pour extraire, de manière non consciente, des régularités statistiques et des probabilités de transition entre les éléments (ex: la probabilité qu'un déterminant soit suivi d'un nom).\cite{II-2-2-ref7}\item \textbf{Nature Implicite et Ineffable :} Les connaissances procédurales sont impénétrables à la conscience. Un locuteur natif utilise le subjonctif correctement sans nécessairement pouvoir expliquer pourquoi. C'est un \og savoir-faire\fg{} (knowledge-how) plutôt qu'un \og savoir-que\fg{} (knowledge-that).\item \textbf{Traitement Séquentiel et Prédictif :} Ce système est spécialisé dans le traitement de séquences en temps réel, qu'il s'agisse de séquences motrices (danser, conduire) ou cognitives (structurer une phrase). Il permet une prédiction rapide de l'élément suivant dans une chaîne, facilitant ainsi la fluidité.\cite{II-2-2-ref2}\item \textbf{Automaticité :} Une fois consolidée, une compétence procédurale s'exécute rapidement, avec une faible variabilité et, surtout, avec un coût attentionnel minimal, libérant ainsi la mémoire de travail pour d'autres tâches (comme la gestion du sens).\end{itemize}

\subsection{Substrats neuroanatomiques}

Le système procédural est enraciné dans les circuits fronto-striataux. Les structures clés
comprennent les ganglions de la base (noyaux gris centraux), en particulier le noyau caudé et le
putamen (le striatum), qui reçoivent des projections du cortex frontal (y compris l'aire de Broca,
impliquée dans la syntaxe) et du cortex pariétal. Le cervelet joue également un rôle crucial dans le
séquençage temporel fin et l'optimisation des automatismes.\cite{II-2-2-ref4}

Ces structures sont dopaminergiques : la dopamine joue un rôle essentiel dans le renforcement des
connexions synaptiques au sein du striatum lors de l'apprentissage par renforcement (réussite de la
communication) ou par prédiction d'erreur.\cite{II-2-2-ref8}



\subsection{Tableau comparatif des systèmes de mémoire dans le langage}

\begin{table}[h!]\small\centering\begin{tabular}{| p{2.3cm} | p{3.5cm} | p{3.5cm} |}\hline\textbf{Dimension} & \textbf{Mémoire Déclarative (Lexique/Règles Explicites)} & \textbf{Mémoire Procédurale (Grammaire/Automatismes)} \\ \hline\textbf{Type de Connaissance} & Explicite, consciente, verbalisable (\og Savoir que\fg{}) & Implicite, inconsciente, ineffable (\og Savoir faire\fg{}) \\ \hline\textbf{Vitesse d'Acquisition} & Rapide (parfois immédiate) & Lente (nécessite pratique massive et répétition) \\ \hline\textbf{Vitesse de Traitement} & Lente (recherche consciente, délibération) & Rapide (millisecondes, réflexe) \\ \hline\textbf{Substrat Anatomique} & Lobe Temporal Médial (Hippocampe), Néocortex & Ganglions de la base (Striatum), Cortex Frontal (Broca) \\ \hline\textbf{Contenu Linguistique} & Mots, irrégularités, règles apprises par cœur & Syntaxe, morphologie régulière, phonologie \\ \hline\textbf{Sensibilité Attentionnelle} & Forte (sature la mémoire de travail) & Faible (fonctionne en \og tâche de fond\fg{}) \\ \hline\textbf{Dépendance à l'Âge} & Stable, s'améliore parfois avec l'âge & Décline potentiellement après l'adolescence (discuté) \\ \hline\textbf{Impact Pédagogique} & Cible de l'enseignement traditionnel (GTM) & Cible de l'immersion et de la pratique communicative \\ \hline\end{tabular}\end{table}Sources : 2

L'un des apports les plus critiques de la neurobiologie à la compréhension de l'impasse pédagogique
est la découverte que ces deux systèmes de mémoire ne fonctionnent pas simplement en parallèle ; ils
interagissent souvent de manière compétitive. Ce phénomène, qualifié d'effet de bascule ou \og see-
saw effect\fg{} par Poldrack, Packard et Ullman, suggère que l'activation intense d'un système peut
physiologiquement inhiber l'autre.\cite{II-2-2-ref4}



\subsection{Mécanismes neurobiologiques de l'inhibition}

Les recherches menées sur des modèles animaux et humains ont mis en évidence une relation
antagoniste entre l'hippocampe (mémoire déclarative) et le noyau caudé (mémoire procédurale).

\begin{itemize}\item \textbf{Inhibition Anatomique :} Il existe des projections inhibitrices directes et indirectes entre les structures du lobe temporal médial et les ganglions de la base. Lorsqu'une tâche d'apprentissage est abordée sous un angle déclaratif (chercher à comprendre et mémoriser consciemment la solution), l'activité accrue dans l'hippocampe envoie des signaux qui diminuent l'activité dans le striatum.\cite{II-2-2-ref12}\item \textbf{Compétition pour les Ressources :} Le cerveau semble optimiser l'apprentissage en sélectionnant le système le plus \og fiable\fg{} pour la tâche en cours. Dans un contexte scolaire où l'erreur est pénalisée et où l'analyse est valorisée, le cerveau privilégie la stratégie déclarative (\og je vais réfléchir à la règle pour ne pas me tromper\fg{}) au détriment de l'essai-erreur nécessaire à l'calibrage procédural.\item \textbf{Modulation Neurochimique :} Des études suggèrent que les niveaux d'œstrogènes peuvent influencer cette balance. Des niveaux élevés d'œstrogènes tendent à favoriser la mémoire déclarative et à inhiber la mémoire procédurale, ce qui a des implications potentielles pour les différences individuelles et de genre dans les stratégies d'apprentissage des langues.\cite{II-2-2-ref4}\end{itemize}

\subsection{L'enseignement explicite comme facteur inhibiteur}

Appliqué à la salle de classe, l'effet de bascule offre une explication biologique à la \og
paralysie par l'analyse\fg{}. Lorsqu'un enseignant demande à un élève de \og bien réfléchir à sa
grammaire\fg{} avant de parler, il force l'activation du système déclaratif.

\begin{itemize}\item \textbf{Conséquence Immédiate :} L'élève accède à la règle stockée dans l'hippocampe/néocortex.\item \textbf{Coût Caché :} Cette activation déclarative \textit{bloque} les processus procéduraux qui auraient permis une production fluide. L'élève se retrouve à \og construire\fg{} sa phrase comme une équation mathématique plutôt que de la \og générer\fg{} comme un flux linguistique.\cite{II-2-2-ref11}\item \textbf{Cercle Vicieux :} Plus l'élève s'appuie sur des règles explicites pour compenser son manque de fluidité, plus il renforce la dominance déclarative et inhibe le développement procédural. Il devient un expert de la règle, mais un novice de l'usage. Poldrack et Packard ont démontré que même si la distraction (qui empêche l'analyse consciente) peut réduire la performance immédiate, elle peut paradoxalement favoriser l'apprentissage procédural à long terme en \og désactivant\fg{} la compétition déclarative.\cite{II-2-2-ref15}\end{itemize}

\subsection{Preuves empiriques dans l'acquisition l2}

Des études utilisant les potentiels évoqués (ERP) et l'imagerie fonctionnelle (fMRI) montrent que
les apprenants L2 exposés à un enseignement explicite (\og voici la règle\fg{}) continuent de
traiter la syntaxe via des mécanismes déclaratifs (signatures N400, typiques du traitement
sémantique/lexical) même après une pratique modérée. À l'inverse, les apprenants exposés à des
conditions implicites (immersion sans règles) développent plus rapidement des signatures neuronales
proches des natifs (P600/LAN, typiques du traitement syntaxique automatique).\cite{II-2-2-ref16}

Cela confirme que l'instruction explicite ne se contente pas de fournir une connaissance ; elle
modifie la trajectoire neurocognitive de l'apprentissage, orientant le cerveau vers une voie qui
n'est pas celle de la fluidité naturelle.

Si l'inhibition procédurale constitue le premier volet de l'impasse, la saturation de la capacité de
traitement en constitue le second. La Théorie de la Charge Cognitive (CLT), développée par John
Sweller, fournit un cadre rigoureux pour analyser pourquoi l'enseignement explicite, en surchargeant
la mémoire de travail, empêche l'apprentissage efficace.\cite{II-2-2-ref19}



\subsection{L'architecture cognitive et le goulot d'étranglement}

La CLT repose sur une caractéristique fondamentale de l'architecture cognitive humaine : la
\textbf{mémoire de travail} est extrêmement limitée.

\begin{itemize}\item \textbf{Capacité :} Elle ne peut traiter qu'environ 4 à 7 éléments d'information nouveaux simultanément.\cite{II-2-2-ref20}\item \textbf{Durée :} Les informations non répétées ou non intégrées disparaissent en quelques secondes (15-30 secondes).\item \textbf{Interaction :} Pour qu'un apprentissage (transfert vers la mémoire à long terme) ait lieu, la mémoire de travail ne doit pas être saturée.\end{itemize}Dans l'apprentissage des langues, Sweller distingue trois types de charges 21 :

1. \textbf{Charge Intrinsèque :} Liée à la complexité de la langue elle-même (vocabulaire inconnu,
structures syntaxiques complexes).

2. \textbf{Charge Extrinsèque :} Liée à la manière dont l'information est présentée (méthode
pédagogique, explications confuses, nécessité de chercher dans un dictionnaire).

3. \textbf{Charge Germane :} La charge cognitive utile consacrée à la création de schémas mentaux.



\subsection{L'anatomie de la saturation en classe de langue}

L'enseignement explicite, en particulier la méthode Grammaire-Traduction, impose une charge
extrinsèque dévastatrice qui sature systématiquement la mémoire de travail.



\subsection{Le coût du "monitor" et du traitement conscient}

Selon l'hypothèse du Monitor de Krashen, l'utilisation consciente de règles pour vérifier la
production (ex: accord sujet-verbe, choix du temps) consomme une quantité massive de ressources
attentionnelles.\cite{II-2-2-ref23} Lorsqu'un élève doit produire une phrase, il doit gérer
simultanément :

\begin{itemize}\item Le plan conceptuel (ce qu'il veut dire).\item La récupération lexicale (trouver les mots).\item L'application algorithmique des règles de grammaire (charge déclarative).\item La gestion phonologique (prononciation).\end{itemize}Cette \og interactivité des éléments\fg{} dépasse largement la capacité de la mémoire de
travail.\cite{II-2-2-ref25} L'élève entre alors en surcharge cognitive : il bégaie, perd le fil de
sa pensée, ou abandonne la complexité grammaticale pour revenir à des structures simplistes.



\subsection{L'effet de l'attention partagée (split-attention effect)}

Un mécanisme spécifique de saturation identifié par Sweller est l'effet d'attention partagée. Cela
se produit lorsque l'apprenant doit diviser son attention entre deux sources d'information
physiquement ou mentalement séparées pour comprendre le tout.\cite{II-2-2-ref26}

\begin{itemize}\item \textbf{Exemple Typique :} Un élève lit un texte en L2 et doit constamment consulter une liste de vocabulaire ou un dictionnaire pour comprendre les mots inconnus. Chaque consultation interrompt le processus de construction du sens et force une réinitialisation partielle de la mémoire de travail. L'intégration mentale ne peut se faire, car les ressources sont épuisées par la navigation entre le texte et l'outil de référence.\cite{II-2-2-ref29}\item \textbf{Traduction Mentale :} La stratégie de traduction (penser en L1 \-\> traduire en L2) est une forme interne d'attention partagée. Elle oblige le cerveau à maintenir deux structures linguistiques parallèles et à effectuer des opérations de conversion coûteuses (mapping), doublant la charge cognitive par rapport au traitement direct en L2.\cite{II-2-2-ref30}\end{itemize}

\subsection{L'effet de redondance}

L'enseignement explicite présente souvent la même information sous plusieurs formes simultanées (ex:
l'enseignant lit une règle qui est projetée au tableau). Contrairement à l'intuition, cela peut
augmenter la charge cognitive si l'élève tente de traiter les deux flux comme des informations
distinctes avant de réaliser qu'ils sont identiques, gaspillant ainsi des ressources de traitement
précieuses.\cite{II-2-2-ref32}



\subsection{Le seuil de vocabulaire et l'effondrement de la compréhension}

Un facteur aggravant majeur de la saturation cognitive est le manque de vocabulaire automatisé. Les
travaux de Paul Nation et Hu ont établi des seuils critiques de couverture lexicale pour la
compréhension.

\begin{itemize}\item \textbf{Le Seuil de 98\% :} Pour qu'un lecteur puisse comprendre un texte de manière autonome et apprendre du vocabulaire par le contexte (incidental learning), il doit connaître déjà 98\% des mots du texte.\cite{II-2-2-ref33}\item \textbf{En dessous de 95\% :} La compréhension est gravement compromise. Le lecteur ne \og lit\fg{} plus ; il \og déchiffre\fg{}. Chaque mot inconnu devient un problème de résolution (problem-solving) qui charge la mémoire de travail.\item \textbf{Conséquence pour la Grammaire :} Si l'élève consacre toute sa mémoire de travail à décoder le sens des mots (processus déclaratif lent), il ne reste \textit{aucune} ressource cognitive pour traiter la structure grammaticale (processus procédural). Le cerveau ignore la syntaxe pour sauver le sens. Ainsi, tant que le vocabulaire n'est pas automatisé, l'entraînement procédural de la grammaire est cognitivement impossible lors de la lecture ou de l'écoute.\cite{II-2-2-ref36}\end{itemize}Face à la saturation du système déclaratif, les pédagogies explicites (comme la méthode PPP :
Présentation, Pratique, Production) promettent une transition : la connaissance déclarative se
transformera, à force d'exercices, en compétence procédurale. Cette vision s'appuie sur la Théorie
de l'Acquisition des Compétences (Skill Acquisition Theory) de Robert DeKeyser, dérivée du modèle
ACT-R d'Anderson.\cite{II-2-2-ref7} Cependant, l'analyse des contraintes temporelles et pédagogiques
révèle que cette conversion est, dans la plupart des contextes scolaires, un mythe.



\subsection{Le modèle act-r et l'exigence de pratique massive}

Selon le modèle ACT-R, la transformation d'une connaissance déclarative (\og savoir que\fg{}) en une
règle de production procédurale (\og savoir faire\fg{}) nécessite trois étapes :

1. \textbf{Stade Cognitif (Déclaratif) :} L'apprenant mémorise la règle.

2. \textbf{Stade Associatif (Procéduralisation) :} L'apprenant applique la règle, détecte les
erreurs, et commence à former des productions spécifiques.

3. \textbf{Stade Autonome (Automatisme) :} La compétence devient rapide et inconsciente.

Le problème réside dans la condition \textit{sine qua non} de ce passage : la \textbf{quantité de
pratique}. La procéduralisation exige une répétition massive, contextuelle et
variée.\cite{II-2-2-ref38} Les estimations pour atteindre l'automaticité dans une compétence
complexe comme la langue se chiffrent en milliers d'heures (souvent citées autour de 600 à 1000
heures pour une compétence fonctionnelle de base, et bien plus pour une fluidité
avancée).\cite{II-2-2-ref41}



\subsection{Le "fossé de la pratique" (the practice gap)}

Dans un cadre scolaire typique (collège/lycée), un élève reçoit environ 2 à 4 heures d'instruction
par semaine. Sur une année scolaire de 30 semaines, cela représente 60 à 120 heures.

\begin{itemize}\item \textbf{Ratio Déclaratif/Procédural en Classe :} Dans une approche traditionnelle, une grande partie de ce temps est consacrée à l'écoute de l'enseignant (explications métalinguistiques), à la lecture de règles et à des exercices de remplissage (drills mécaniques). Ces activités maintiennent l'élève dans le stade cognitif/déclaratif.\item \textbf{Temps de Pratique Réelle :} Le temps de parole effectif (production active) par élève dans une classe de 30 personnes se mesure souvent en secondes ou en minutes par cours.\item \textbf{Impossibilité Mathématique :} Il est impossible d'atteindre le volume de pratique nécessaire à la procéduralisation (loi de puissance de la pratique) avec une densité de pratique aussi faible. Les règles explicites s'accumulent (saturation déclarative), mais ne sont jamais pratiquées assez intensément pour basculer vers le système procédural.\cite{II-2-2-ref43} L'élève reste \og coincé\fg{} au stade 1 de DeKeyser.\end{itemize}

\subsection{Le débat de l'interface : krashen vs dekeyser}

La question de savoir si la connaissance explicite \textit{peut} devenir implicite est au cœur du
débat sur l'interface.

\begin{itemize}\item \textbf{Position de l'Interface Forte (DeKeyser) :} Oui, avec assez de pratique. Mais comme démontré, \og assez de pratique\fg{} est irréalisable scolairement.\cite{II-2-2-ref45}\item \textbf{Position de la Non-Interface (Krashen) :} Non, les systèmes sont dissociés. L'apprentissage (learning) ne devient jamais acquisition. La connaissance consciente sert uniquement de Moniteur pour l'édition, mais ne génère pas la parole spontanée.\cite{II-2-2-ref23}\item \textbf{Perspective Neurocognitive (Ullman) :} Le modèle DP suggère une position intermédiaire nuancée mais pessimiste pour l'école. Bien que des connexions existent, le transfert est biologiquement coûteux et inefficace. Le cerveau préfère développer le système procédural directement par l'exposition (immersion) plutôt que par le détour déclaratif. Tenter de forcer ce détour via l'enseignement explicite revient à nager à contre-courant neurobiologique.\cite{II-2-2-ref2}\end{itemize}

\subsection{L'échec de la méthode grammaire-traduction (gtm)}

La GTM illustre parfaitement l'impasse de la non-conversion. En traitant la langue comme un objet
d'analyse intellectuelle (traduction de textes littéraires, dissection syntaxique), elle forme des
experts en linguistique descriptive, mais des locuteurs dysfonctionnels.

\begin{itemize}\item \textbf{Absence de Théorie :} Comme le soulignent Richards et Rodgers, la GTM est une méthode \og sans théorie\fg{}, née de la tradition de l'enseignement du latin (langue morte, donc sans besoin procédural de conversation) et appliquée par inertie aux langues vivantes.\cite{II-2-2-ref47}\item \textbf{Conséquences :} Elle produit des apprenants capables d'analyser une phrase complexe à l'écrit (compétence déclarative élevée) mais incapables de commander un repas sans hésitation (compétence procédurale nulle). Le savoir est \og inerte\fg{}.\cite{II-2-2-ref48}\end{itemize}L'impasse de l'enseignement explicite n'est pas un concept abstrait, mais une réalité physiologique
observable. Elle résulte de la convergence de trois mécanismes d'échec simultanés :

1. \textbf{Saturation de la Mémoire de Travail (Cognitive Overload) :} L'exigence de manipuler
consciemment des règles, du vocabulaire et du sens sature la capacité de traitement (4-7 éléments),
provoquant l'effondrement de la communication et l'arrêt de l'apprentissage (pas de consolidation).

2. \textbf{Inhibition Procédurale (See-Saw Effect) :} L'effort conscient d'analyse active
l'hippocampe, qui envoie des signaux inhibiteurs aux ganglions de la base, bloquant
physiologiquement les mécanismes d'automatisation nécessaires à la fluidité.

3. \textbf{Carence de Pratique (Practice Gap) :} Le temps scolaire est cannibalisé par l'instruction
explicite (déclarative), ne laissant qu'une fraction marginale pour la pratique communicative,
rendant la conversion procédurale (même si elle était possible) mathématiquement inatteignable.



\subsection{Tableau récapitulatif : les mécanismes de l'impasse}

\begin{table}[h!]\small\centering\begin{tabular}{| p{2cm} | p{2.2cm} | p{2.2cm} | p{2.2cm} |}\hline\textbf{Mécanisme} & \textbf{Processus Pédagogique Traditionnel} & \textbf{Conséquence Neurocognitive} & \textbf{Résultat Observable} \\ \hline\textbf{Biais Déclaratif} & Enseignement de règles explicites (\og Focus on Forms\fg{}) & Activation de l'hippocampe / Cortex préfrontal & \og Savoir\fg{} la règle mais ne pas pouvoir l'utiliser. \\ \hline\textbf{Effet de Bascule} & Focalisation sur l'analyse et la précision (Accuracy) & Inhibition du Striatum (Ganglions de la base) & \og Paralysie par l'analyse\fg{}, hésitation, perte de fluidité. \\ \hline\textbf{Attention Partagée} & Traduction mentale, recours au dictionnaire & Saturation de la Mémoire de Travail & Fatigue cognitive rapide, incompréhension globale, décrochage. \\ \hline\textbf{Manque d'Input} & Priorité à la production précoce (Output) et aux exercices & Carence statistique pour le système procédural & Absence d'intuition grammaticale (\og Sprachgefühl\fg{}), erreurs fossilisées. \\ \hline\textbf{Seuil Lexical} & Listes de mots hors contexte, vocabulaire faible & Absence de reconnaissance automatique des mots & Le traitement grammatical est sacrifié pour le décodage lexical. \\ \hline\end{tabular}\end{table}Sources : 11

Si l'enseignement explicite mène à une impasse, quelles sont les voies de sortie suggérées par la
recherche? Les données pointent vers une réorientation radicale des priorités pédagogiques.



\subsection{Réduire la charge déclarative}

Il est essentiel de minimiser l'explication métalinguistique, surtout aux niveaux débutants et
intermédiaires. L'enseignant doit résister à la tentation d'expliquer \og pourquoi\fg{} une phrase
est correcte, car cette explication active le système inhibiteur. L'objectif est de réduire la
charge extrinsèque pour libérer la mémoire de travail.\cite{II-2-2-ref22}



\subsection{Priorité à l'input compréhensible}

Le système procédural étant un système d'apprentissage statistique, il a besoin de données.
L'approche de l'immersion (\og Input Flood\fg{}) et les méthodes basées sur la compréhension
(Comprehensible Input) fournissent au cerveau la matière première nécessaire pour extraire les
régularités grammaticales sans passer par le filtre conscient.\cite{II-2-2-ref23} Le vocabulaire
doit être acquis en contexte pour faciliter l'accès lexical rapide, condition préalable à
l'automatisation syntaxique.\cite{II-2-2-ref37}



\subsection{Automatisation par la pratique communicative}

La pratique ne doit pas être une répétition mécanique de règles (drills), mais une interaction
signifiante où l'attention est portée sur le message. Cela favorise l'engagement des circuits
dopaminergiques du striatum, essentiels à la consolidation procédurale.\cite{II-2-2-ref51}

En conclusion, l'expression \og l'impasse de l'enseignement explicite\fg{} décrit avec précision une
réalité neurobiologique : on ne peut pas forcer le cerveau à automatiser une compétence motrice
complexe (la parole) en surchargeant son système de mémoire encyclopédique (les faits). L'avenir de
l'enseignement des langues réside dans des méthodes qui respectent la dissociation de ces systèmes
et qui cessent de traiter la langue comme une connaissance académique pour la traiter comme ce
qu'elle est vraiment : une habileté procédurale vivante.

